\chapter{Configurações Vencedoras da Otimização Conjunta}
\label{ap:configuracoes}

A \autoref{tab:configuracoesVencedoras} lista, para cada cena e segmentador, a configuração escolhida pela otimização bayesiana com processo gaussiano (\autoref{sec:desHPO}) no experimento 03: 16 \textit{trials} por segmentador, os quatro primeiros sorteados, treino completo (85\% das amostras), teto de 200 épocas, parada após 10 épocas sem que a acurácia de validação melhore 0,001 e seleção pela acurácia de validação. A coluna da configuração do segmentador indica qual das três configurações do espaço de busca foi escolhida: a padrão, que no SEEDS é a do original e no LSC e no Watershed gera o mesmo número de regiões que o SEEDS, ou uma das duas variações em torno dela (\autoref{qua:espaco}). A configuração vencedora foi retreinada com 5 sementes sob as mesmas regras; como treino e parada são os da linha de base, a OA de teste é comparável à da \autoref{tab:linhaBase}, que é de uma única semente.

Em Indian Pines, nos dois segmentadores já medidos, a busca manteve a configuração padrão do segmentador e escolheu os mesmos hiperparâmetros da CNN: taxa de aprendizado de 0,0003, \textit{dropout} de 0,2, sem \textit{weight decay}, lote 64 e 15 componentes do PCA, metade dos 30 da configuração original. A melhor validação da busca ficou 1,6 ponto acima da OA média de teste no SEEDS e 2,4 no LSC, o efeito esperado de escolher o melhor de 16 \textit{trials} pela própria validação. O desvio de OA entre as cinco sementes ficou abaixo de 0,7 ponto nos dois. \pendente{experimento 03: linha do Watershed em Indian Pines e linhas de University of Pavia e Salinas Valley; se alguma configuração tiver desvio alto entre sementes, descrever as sementes em que o treino colapsa}

A coluna Épocas traz a média de épocas treinadas no retreino. Sem meta de validação, o treino para após 10 épocas sem melhora ou no teto de 200; em Indian Pines, as médias foram de 69,6 épocas no SEEDS e 82,6 no LSC, acima das 60 do notebook original. \pendente{experimento 03: médias de épocas das demais linhas}

\begin{table}[H]
    \centering
    \caption{Configurações vencedoras por cena e segmentador (experimento 03; 5 sementes, treino completo)}
    \footnotesize
    \setlength{\tabcolsep}{1.2pt}
    \begin{tabular}{|l|c|c|C{0.9cm}|c|c|l|c|c|c|}
        \hline
        \rowcolor[HTML]{C0C0C0}
        \textbf{Segmentador} & \textbf{Taxa} & \textbf{\textit{Dropout}} & \textbf{\textit{Weight decay}} & \textbf{Lote} & \shortstack{\textbf{Comp.}\\\textbf{do PCA}} & \shortstack{\textbf{Configuração do}\\\textbf{segmentador}} & \shortstack{\textbf{Validação}\\\textbf{da busca}} & \textbf{Épocas} & \textbf{OA de teste} \\
        \hline
        \multicolumn{10}{|l|}{\cellcolor[HTML]{EFEFEF}\textbf{Indian Pines}} \\
        \hline
        SEEDS & 0,0003 & 0,20 & 0 & 64 & 15 & 350 sp., 200 it., \textit{prior} 2 & 0,9480 & 69,6 & 0,9324\,$\pm$\,0,0062 \\
        \hline
        LSC & 0,0003 & 0,20 & 0 & 64 & 15 & 16 px; razão 0,075 & 0,9389 & 82,6 & 0,9152\,$\pm$\,0,0045 \\
        \hline
        Watershed & -- & -- & -- & -- & -- & -- & -- & -- & -- \\
        \hline
        \multicolumn{10}{|l|}{\cellcolor[HTML]{EFEFEF}\textbf{University of Pavia}} \\
        \hline
        SEEDS & -- & -- & -- & -- & -- & -- & -- & -- & -- \\
        \hline
        LSC & -- & -- & -- & -- & -- & -- & -- & -- & -- \\
        \hline
        Watershed & -- & -- & -- & -- & -- & -- & -- & -- & -- \\
        \hline
        \multicolumn{10}{|l|}{\cellcolor[HTML]{EFEFEF}\textbf{Salinas Valley}} \\
        \hline
        SEEDS & -- & -- & -- & -- & -- & -- & -- & -- & -- \\
        \hline
        LSC & -- & -- & -- & -- & -- & -- & -- & -- & -- \\
        \hline
        Watershed & -- & -- & -- & -- & -- & -- & -- & -- & -- \\
        \hline
    \end{tabular}
    \captionsetup{justification=raggedright}
    \caption*{Fonte: Elaborada pelo autor. Taxa: taxa de aprendizado. Configuração do segmentador: no SEEDS, superpixels pedidos (sp.), iterações (it.) e \textit{prior}; no LSC, lado da região em pixels e razão; no Watershed, o número de marcadores, dispostos em grade. Comp. do PCA: 15, 30 ou 45 em Indian Pines e 8, 15 ou 22 nas outras cenas. Validação da busca: melhor acurácia de validação entre os \textit{trials}; épocas: média das épocas treinadas no retreino; OA de teste: média e desvio padrão entre as 5 sementes.}
    \label{tab:configuracoesVencedoras}
\end{table}
